\documentclass[12pt, titlepage]{article}
\usepackage{hyperref}

\title{Feedback Controls Lab Weekly Report \\ \normalsize{Week 5}}
\author{Aydan Vissing, Brady Schick, Gabriel Southwell, Simon Ross}
\date{\today}

\begin{document}
\maketitle

\section*{Aydan Vissing - CFO (9 hrs spent)}

I wasn't able to put quite as much work in this week since I was sick. However, I was able to mostly assemble the rig, there being a piece I have to reprint since it broke during the assembly process, but that won't be long. I did research on how to use QT design studio rather than Designer and am continuing development on the front end application. 

\section*{Brady Schick - CTO (12 hrs spent)}

I pretty much just spent more time reading through Crazyflie code. I expect to update the simulation to line up with how the Crazyflie code will work and possibly help debug the current control code for the drone this week if things are ready to start doing that.

\section*{Gabriel Southwell - CSO (11.5 hrs spent)}

This week I finished swapping out all the critical components that I know of. 

The ESP\_Drone project used a separate eeprom on the PCB to store and load configuration and tuning values. Since we don't have said eeprom, I replaced it with storing the configs on the 1k NVS partition that we were already using for WiFi configs. 

I also took the time to replace the buzzer component with the one I had already written. That was a fairly easy and quick fix. 

I am now reasonably sure that all the mission critical stuff can work. I spent some time at the end of the week learning about how to use the cflib python library to communicate with the drone. If it works correctly, I will be able to flash it and connect to the WiFi and see what systems have booted correctly. If everything important works, then I can use it to sent some motor controls or check data logging. 

Something that will big hurdle in the nearish future is figuring out what to do with the time of flight sensors. The ESP\_Drone does have functionality for vl53l1s, but it only one and it is fed into the sensor fusion to improve state estimation. 

\section*{Simon Ross - CEO (12 hrs spent)}

This week, I found a bug in the battery checking code and finished the working implementation. I'll have to do a little more soldering to get the voltage divider values correct and implement logging code and then I will move on to scripting. Sounds like Gabe finished the flight code, so we will likely be doing flight testing this week as well. I need to work on completing work more quick-and-dirty than I have been, which would speed up my work---as mentioned, it doesn't need to be as pretty as I'm making it.

\end{document}